\chapter{生产级推理引擎（vLLM / SGLang / TRT-LLM）压测实践}

\section{毕业设计：工业级推理引擎实践（vLLM / llama.cpp）}

\begin{mathBox}[核心数学定理推导]
这是整条路线的\textbf{毕业设计}。前面的 Lab 03\textasciitilde{}13 里你已经亲手实现了推理引擎的每一个核心部件；现在去使用真正的引擎，并且\textbf{每一个观测到的数字，都要能用前面的公式解释}。

前置：第 8、10、11、12 篇；Lab 08（实测带宽）、Lab 07b（手写 Qwen）、Lab 12b（PagedAttention）。
\end{mathBox}

\section{原则：先预测，再测量，再解释}

每做一个实验，先用 \texttt{lab08\_memory\_and\_roofline/memory\_calculator.py} 和第 8 篇的公式\textbf{写下你的预测}，再运行，最后解释差距。这比跑出一个漂亮的数字重要得多。

\vspace{0.8em}\hrule\vspace{0.8em}

\section{[测试] 任务 1：用 vLLM 部署，并测出吞吐-延迟曲线}

\textbf{安装}：建议单独建一个虚拟环境（vLLM 会固定自己需要的 PyTorch 版本，可能和本仓库的 \texttt{.venv} 冲突）。RTX 50 系显卡需要支持 CUDA 12.8 的较新版本，按 \href{https://docs.vllm.ai/}{vLLM 官方安装文档} 操作。WSL2 下可以正常使用。

\begin{lstlisting}[language=bash]
vllm serve Qwen/Qwen2.5-1.5B-Instruct --max-model-len 4096 --gpu-memory-utilization 0.85
\end{lstlisting}

\textbf{先预测}：
\begin{itemize}
  \item Qwen2.5-1.5B：28 层、2 个 KV 头、head\_dim 128 $\rightarrow$ 每 token KV = $2 \times 28 \times 2 \times 128 \times 2 = 28{,}672$ 字节 $\approx$ 28 KB；
  \item 权重 BF16 约 3.1 GB，按你实测的带宽（本机空闲热态约 380 GB/s，负载下会更低），单流 Decode 上限 $\approx$ 120 token/s；
  \item 16GB $\times$ 0.85 扣掉权重和激活后，大约还剩多少 GB 给 KV？能放多少 token？
\end{itemize}

\textbf{再测量}：
\begin{enumerate}
  \item 看 vLLM 启动日志里的 KV Cache 容量（不同版本措辞不同，通常会打印 KV Cache 的 token 数或 GPU block 数与最大并发）。和你的预测一致吗？block 数 $\times$ 16 = token 数？
  \item 压测：
\end{enumerate}

\begin{lstlisting}[language=bash]
   python3 capstone_engines_practice/bench_openai_server.py --concurrency 1 4 16 64 --max-tokens 256
   
\end{lstlisting}

\begin{enumerate}
  \item 画出（或列出）"并发 $\rightarrow$ 总吞吐、TPOT P50"的关系。
\end{enumerate}

\begin{warningBox}[避坑指南与核心警示]
{}[注意] \textbf{压测本身有三个坑，不知道就会得出错误的结论}（脚本已在文件头 docstring 里写明，这里再强调一次）：
- \textbf{前缀缓存会污染跨档位的对比}：脚本各档位复用同一批问题，档位之间不清缓存，而 vLLM 默认开启前缀缓存 $\rightarrow$ 靠后的档位 TTFT 偏低。要得到干净的曲线，请在档位之间重启服务。
- \textbf{输出长度必须一致}：默认加了 \texttt{ignore\_eos: true}（vLLM/SGLang 扩展），让每个请求都恰好输出 \texttt{--max-tokens} 个 token。否则"总吞吐 tok/s"测的是模型有多话唠，不是引擎有多快。标准 OpenAI 服务会忽略这个字段，那就只能靠 \texttt{max\_tokens} 近似。
- \textbf{必须预热}：脚本每档先发 \texttt{--warmup}（默认 2）个不计入统计的请求。GPU 从空闲频率爬到稳态要时间——lab08 里同一份代码在不同时钟状态下读到过 217\textasciitilde{}380 GB/s，差近一倍。第一档不预热会显著偏慢。
\end{warningBox}

\begin{noteBox}[核心导读与背景]
{}[重点] 另外，脚本会\textbf{自检} \texttt{usage.completion\_tokens} 和实际收到的内容分片数是否一致：不一致时说明服务端一个 SSE chunk 带了多个 token（llama.cpp 常见），此时 TPOT 按分片间隔计算并打印告警——否则 TPOT 会被系统性高估。样本不足（例如 \texttt{--max-tokens 1}）时它会打印"样本不足"而不是崩溃。
\end{noteBox}

\textbf{要能解释}：
\begin{itemize}
  \item 并发 1 的单流速度为什么低于 120 token/s，但比 Lab 07b 的手写实现（约 20ms/token）快得多？（提示：CUDA Graph、融合 kernel，第 12 篇 §7）
  \item 并发从 1 到 16，总吞吐为什么几乎线性增长、TPOT 却几乎不变？到多少并发开始 TPOT 明显变差？（第 8 篇 §4：拐点约 130）
\end{itemize}

\vspace{0.8em}\hrule\vspace{0.8em}

\section{[测试] 任务 2：前缀缓存}

\begin{lstlisting}[language=bash]
python3 capstone_engines_practice/bench_openai_server.py --shared-prefix-tokens 3000 --concurrency 4 --max-tokens 64
\end{lstlisting}

分别在开启和关闭前缀缓存时运行（新版 vLLM 默认开启；关闭参数通常是 \texttt{--no-enable-prefix-caching}，以你的版本 \texttt{vllm serve --help} 为准）。

\textbf{预测}：3000 token 的 Prefill 按第 8 篇 $2N \cdot S / (\pi \cdot \text{MFU})$ 估算要多少毫秒？命中缓存后 TTFT 应该降到多少？

\vspace{0.8em}\hrule\vspace{0.8em}

\section{[测试] 任务 3：长 Prefill 对 Decode 的干扰（Chunked Prefill）}

同时跑两个压测：一个并发 16、短问题、长输出（模拟正在聊天的用户）；另一个偶尔发送 3000+ token 的长提示词。观察第一个的 \textbf{TPOT P99} 是否出现尖刺。
调整 \texttt{--max-num-batched-tokens}（每步 token 预算）重复实验，用 Lab 12b 实验 5 的模型解释你看到的现象。

\vspace{0.8em}\hrule\vspace{0.8em}

\section{[测试] 任务 4：llama.cpp + GGUF 量化}

\begin{lstlisting}[language=bash]
git clone https://github.com/ggml-org/llama.cpp && cd llama.cpp
cmake -B build -DGGML_CUDA=ON && cmake --build build --config Release -j
# 从 Hugging Face 上 Qwen 官方的 GGUF 仓库下载 Qwen2.5-7B-Instruct 的 Q4_K_M 版本
./build/bin/llama-bench -m <你的 gguf 文件> -ngl 99
\end{lstlisting}

\texttt{llama-bench} 会报告 \texttt{pp512}（Prefill 512 token 的吞吐）和 \texttt{tg128}（Decode 的吞吐）。

\textbf{预测}：
\begin{itemize}
  \item Q4\_K\_M 平均约 4.8 bit/参数，7.6B 参数约 4.6 GB $\rightarrow$ \texttt{tg} 上限 $\approx$ 380 / 4.6 $\approx$ 80 token/s；
  \item \texttt{pp} 是算力受限：约 $\frac{48 \times 10^{12} \times \text{MFU}}{2 \times 7.6 \times 10^9}$ token/s。
  \item BF16 的 7B（15.2 GB）在 16GB 卡上几乎放不下——这就是你这张卡上量化的第一价值（第 11 篇）。
\end{itemize}

\textbf{要能解释}：为什么量化让 \texttt{tg} 提升了约 3 倍，而 \texttt{pp} 提升远没有那么多，甚至可能变慢？（提示：W4A16 在计算前要反量化，Prefill 是算力受限的）

\vspace{0.8em}\hrule\vspace{0.8em}

\section{[测试] 任务 5：投机解码}

在 vLLM（配置草稿模型或 n-gram 投机）或 llama.cpp（\texttt{llama-speculative}，目标 Qwen2.5-7B + 草稿 Qwen2.5-0.5B）上开启投机解码，测单流速度与接受率。
用第 10 篇的公式 $\frac{1 - \alpha^{K+1}}{(1 - \alpha)(Kc + 1)}$ 代入实测接受率和两模型的耗时比，预测加速比并与实测对比。再把并发调到 32，投机解码还有收益吗？为什么？

\vspace{0.8em}\hrule\vspace{0.8em}

\section{[阅读] 任务 6：源码阅读路线}

读源码的顺序很重要——从小到大：

\begin{enumerate}
  \item \textbf{\href{https://github.com/GeeeekExplorer/nano-vllm}{nano-vllm}}（约 1200 行 Python，实现了 vLLM 的核心：调度、块管理、前缀缓存、CUDA Graph、张量并行）。对照表：
\end{enumerate}

\begin{center}
\small
\begin{tabularx}{\textwidth}{p{3.5cm} X}
\toprule
\textbf{nano-vllm 里的部分} & \textbf{本仓库对应} \\
\midrule
调度器（Prefill / Decode、抢占） & Lab 12a、第 12 篇 §2、§3.4 \\
块管理器（引用计数、哈希前缀缓存） & Lab 12b 的 \texttt{BlockAllocator}、实验 4 \\
模型实现（Qwen） & Lab 07b 的 \texttt{QwenFromScratch} \\
注意力（调用 FlashAttention 的变长 / 分页接口） & Lab 09、第 9 篇 \\
CUDA Graph 捕获 & Lab 07b 性能分析、练习 5 \\
张量并行的线性层 & Lab 13、第 13 篇 \\
\bottomrule
\end{tabularx}
\end{center}

\begin{enumerate}
  \item \textbf{\href{https://github.com/pytorch-labs/gpt-fast}{gpt-fast}}（不到 1000 行，用 \texttt{torch.compile} + 量化 + 投机解码把单流速度逼近带宽极限）。
  \item \textbf{vLLM 本体}：先读官方文档的架构设计部分，再按"入口 $\rightarrow$ 引擎核心 $\rightarrow$ 调度器 $\rightarrow$ KV Cache 管理 $\rightarrow$ 模型执行器 $\rightarrow$ 注意力后端"的顺序读。
  \item \textbf{llama.cpp}：从 \texttt{ggml} 的计算图与量化格式（\texttt{Q4\_K} 等）入手，看它如何在 CPU / GPU 上做融合的反量化矩阵乘。
\end{enumerate}

\vspace{0.8em}\hrule\vspace{0.8em}

\section{[OK] 毕业标准}

完成下面这件事，你就真正走完了这条路线：

\begin{mathBox}[核心数学定理推导]
选定一个模型和你的 5060 Ti，写一份一页纸的报告：\textbf{预测}它在 vLLM 上的单流速度、最大并发、TTFT（带/不带前缀缓存）、W4 量化后的速度，给出每个预测的公式；然后\textbf{实测}，逐项解释预测与实测之间的差距来自哪里。
\end{mathBox}

如果每一个差距你都能说清楚是哪一篇、哪一个公式、哪一个工程因素造成的，你就不再是"库的使用者"了。


\section{实验核心源码精选与解析}

本实验配套有完整可运行的工业级验证代码，重点实现细节如下：


\subsection{源码实现：\texttt{bench\_openai\_server.py}}

\lstinputlisting[language=Python, caption={\texttt{bench\_openai\_server.py}}]{code/capstone_engines_practice_bench_openai_server.py}

